% STATUS: reviewed
\chapter{Įvadas}\label{ch:intro}

Šiame skyriuje paaiškinama, kam skirtas \product{}, ką jis daro ir kuriose jo vietose dirba
kalbos modelis. Skyrius baigiamas sistemos skaičiais, kursu, kuriam ji sukurta, ir tuo, ko
sistema neteigia daranti.
Viskas, kas čia paminėta, paaiškinama vėlesniuose skyriuose; nuorodos parodo, kuriuose.

\section{Problema}\label{sec:intro:problem}

Finansų įstaiga Lietuvoje vienu metu turi laikytis kelių dešimčių teisės aktų: tiesiogiai ją
saistančių Europos Sąjungos reglamentų, per nacionalinę teisę ją pasiekiančių direktyvų,
Lietuvos įstatymų ir priežiūros institucijų nutarimų, nurodymų bei gairių. Pavyzdžiui,
kriptoturto paslaugų teikėjas savo pareigų turi ieškoti kriptoturto rinkų reglamente ir jo
techniniuose standartuose, skaitmeninės veiklos atsparumo reglamente, pinigų plovimo
prevencijos ir lėšų pervedimų taisyklėse, duomenų apsaugos reglamente ir nacionaliniuose
įstatymuose, kurie už visų šių reikalavimų pažeidimus nustato baudas. Už atitiktį atsakingas
asmuo turi nustatyti keturis dalykus ir nustatyti juos iš naujo kaskart, kai pasikeičia
įstatymas ar pati įmonė:

\begin{enumerate}
  \item \textbf{Kurie teisės aktai taikomi} šiai įmonei, atsižvelgiant į tai, kas ji yra ir ką
        ji daro.
  \item \textbf{Ką kiekvienas iš jų įpareigoja} įmonę daryti, nuostata po nuostatos.
  \item \textbf{Ar pačios įmonės dokumentai} --- jos politikos ir procedūros --- įgyvendina
        kiekvieną pareigą ir kur jie nepakankami.
  \item \textbf{Ką daryti su spragomis}: kokia tvarka ir kas atsakingas už kiekvieną veiksmą.
\end{enumerate}

Rankiniu būdu tam prireiktų kelių savaičių skaitymo. O skaitymas kaip tik ir yra toks darbas, kurį kalbos
modelis atlieka greitai, bet, paliktas vienas, nepatikimai: jis nurodys straipsnį, kuriame
parašyta kas kita, pacituos sakinį, kurio politikoje nėra, arba pripažins pareigą įvykdyta
vien todėl, kad politikoje paminėta ta tema. \product{} sukurtas aplink šią įtampą. Skaityti
jis leidžia modeliams, o jų iškvietimus apsupa trimis dalykais: paties teisės akto tekstu,
atsisiųstu iš oficialaus šaltinio, kad teisė būtų skaitoma, o ne prisimenama; patikromis,
kurias atlieka įprastas kodas, kad citata ar nuoroda būtų patikrinta, o ne priimta už gryną
pinigą; ir žmogumi, kuris sprendžia visur, kur sprendimas turi pasekmių.

\section{Ką sistema daro}\label{sec:intro:does}

\product{} (taip jis vadinasi savo pirmąja kalba, lietuvių) yra savo serveryje diegiama
žiniatinklio programa. Prisijungęs naudotojas patenka į pokalbį su \term{pokalbių
asistentu}; šalia pokalbio yra puslapiai rezultatams peržiūrėti ir administruoti. Toliau sistemos funkcijos
išvardytos ta tvarka, kuria su jomis susiduria naudotojas.

\begin{description}
  \item[Saugo teisę kaip tekstą.] \term{Teisės aktų bazė} iš 78 kataloge įrašytų aktų, iš kurių
    68 naudojami, atsisiunčiama iš oficialių Europos Sąjungos ir Lietuvos mašininio skaitymo
    šaltinių (dvejos gairės įkeliamos kaip failai), jos struktūra atpažįstama iki straipsnio,
    dalies ir punkto, o tekstai saugomi redakcijomis. Išvados ir atsakymai nurodo šios bazės
    nuostatas ir gali parodyti jų žodžius (\cref{ch:corpus}).
  \item[Paverčia teisę pareigų registru.] Kiekvienas aktas perskaitomas vieną kartą, o tai, ką
    jis įpareigoja daryti, užrašoma kaip \term{pareigos} --- kas ką privalo, kokiomis sąlygomis,
    iki kada ir koks pareigos reikšmingumas --- kiekviena susieta su nuostatomis, iš kurių kyla.
    Tada antrasis modelio iškvietimas kiekvieną pareigą patikrina pagal jos nuostatą, o
    kiekvienos pareigos patvirtinimą sistemos sandara palieka žmogui (\cref{ch:obligations};
    kaip buvo peržiūrėtas kartu su sistema pateikiamas registras, nurodyta
    \cref{ch:control}).
  \item[Susipažįsta su įmone.] Naudotojas pokalbyje aprašo įmonę ir įkelia jos politikas;
    asistentas savo įrankiais įrašo profilį ir dokumentus. Kiekvienas dokumentas indeksuojamas
    paieškai, o jį perskaičius sudaromas struktūrizuotas \term{inventorius} to, ką jis nustato
    (\cref{ch:company}).
  \item[Atsako į klausimus apie teisę.] Pokalbis ieško teisės aktų bazėje, cituoja tikslią
    nuostatą ir tikrina, ar straipsnis, kurį nurodo dokumentas, tikrai reglamentuoja tai, ką
    dokumentas teigia (\cref{ch:chat,ch:company}).
  \item[Rengia atitikties ataskaitą.] Analizė nusprendžia, kurie aktai taikomi įmonei,
    sustoja, kad žmogus patvirtintų šį sąrašą (nebent analizę pradėjęs asmuo nusprendė leisti
    jai vykti be sustojimo), tada kiekvieną tų aktų pareigą, tinkančią įmonei, palygina su
    įmonės dokumentais. Rezultatas --- \term{ataskaita}: kiekvienai tokiai pareigai
    \term{išvada} su būsena, iš dokumentų cituotais įrodymais, rekomendacija ir, jei už spragą
    aktas numato sankcijas, tuo, kiek gali kainuoti pažeidimas; \term{veiksmų planas},
    apimantis visas spragas; ir santrauka valdybai.
    Žmogus peržiūri išvadas ir patvirtina ataskaitą. Ji eksportuojama į PDF ir Excel
    (\cref{ch:report}).
  \item[Seka veiksmus laike.] Patvirtintos ataskaitos planas perkeliamas į įmonės
    \term{veiksmų registrą}, kuriame veiksmas uždaromas, kai vėlesnė patvirtinta ataskaita
    nustato, kad jo pareigos įvykdytos, arba kai žmogus jį uždaro rankiniu būdu su pastaba,
    paaiškinančia kodėl (\cref{ch:register}).
  \item[Rengia atitikties programos dokumentų projektus.] Iš tų pačių dokumentų sistema parengia
    trijų \term{programos dokumentų} projektus: atitikties politiką su dar trūkstamų dokumentų sąrašu;
    valdymo struktūrą, atitikties procesą ir jo atsakomybių matricą; teisinės rizikos
    kriterijus ir registrą. Kodas kiekvieną patikrina pagal tos rūšies dokumento taisykles, o
    žmogus jį redaguoja ir patvirtina (\cref{ch:programme}).
  \item[Rodo, kaip veikia pati.] Administratorius tvarko teisės aktų bazę, peržiūri pareigas,
    skaito ir redaguoja kiekvieną instrukciją, mato kiekvieną užregistruotą modelio
    iškvietimą su jo kaina ir kiekvieną bet kieno atliktą pakeitimą
    (\cref{ch:console,ch:prompts,ch:security}).
\end{description}

\section{Nuo teisės iki ataskaitos viename paveiksle}\label{sec:intro:picture}

\Cref{fig:intro:map} rodo visą sistemą kaip keturias juostas. Pirmoji vykdoma kartą
kiekvienam aktui, gerokai anksčiau, nei atsiranda kokia nors įmonė: tai brangus, kruopštus
teisės skaitymas, kurį pakartotinai panaudoja kiekviena ataskaita. Antroji kiekvieną
dokumentą perskaito vieną kartą; kiekviena ataskaita pakartoja tik jo nuorodų patikrą.
Trečioji juosta --- ataskaita, ketvirtoji --- programos dokumentas. Pokalbių asistentas
stovi virš jų: savo įrankiais jis skaito visas keturias, įrašo įmonės dokumentus ir pradeda
ataskaitą ar programos dokumentą, tačiau niekada pats neatlieka jų darbo.

\begin{figure}[H]
\centering
\begin{tikzpicture}[
    n/.style={text width=15.5mm, inner xsep=3pt, inner ysep=4pt, minimum height=11.5mm, font=\sffamily\scriptsize},
    node distance=5mm and 3.2mm]
  % --- the law
  \node[ext, n] (src) {oficialūs šaltiniai};
  \node[svc, n, right=of src] (corp) {teisės aktų bazė: aktai, nuostatos};
  \node[ai, n, right=of corp] (ext) {išskiriamos pareigos};
  \node[ai, n, right=of ext] (jud) {vertintojas tikrina kiekvieną};
  \node[man, n, right=of jud] (rev) {peržiūri žmogus};
  \node[store, n, right=of rev, text width=15mm] (reg) {pareigų registras};
  \draw[flow] (src) -- (corp); \draw[flow] (corp) -- (ext); \draw[flow] (ext) -- (jud);
  \draw[flow] (jud) -- (rev); \draw[flow] (rev) -- (reg);
  % --- the company
  \node[data, n, below=13mm of src, text width=13mm] (docs) {įmonės dokumentai};
  \node[ai, n, right=of docs] (idx) {indeksavimas: vektoriai};
  \node[ai, n, right=of idx] (inv) {kiekvieno dokumento inventorius};
  \node[ai, n, right=of inv] (cit) {nuorodų patikra};
  \draw[flow] (docs) -- (idx); \draw[flow] (idx) -- (inv); \draw[flow] (inv) -- (cit);
  % --- the report
  \node[ai, n, below=13mm of docs] (app) {kurie aktai taikomi};
  \node[man, n, right=of app] (scope) {apimtis patvirtinta};
  \node[ai, n, right=of scope] (gap) {spragų analizė: išvados};
  \node[chk, n, right=of gap] (chk) {citatos rastos, sankcijos susietos};
  \node[ai, n, right=of chk] (plan) {veiksmų planas, santrauka};
  \node[man, n, right=of plan] (sign) {peržiūrėta, patvirtinta};
  \node[chk, n, right=of sign] (areg) {veiksmų registras};
  \draw[flow] (app) -- (scope); \draw[flow] (scope) -- (gap); \draw[flow] (gap) -- (chk);
  \draw[flow] (chk) -- (plan); \draw[flow] (plan) -- (sign); \draw[flow] (sign) -- (areg);
  % --- work products
  \node[ai, n, below=13mm of app] (wp) {parengiami politika, struktūra, rizikos registras};
  \node[chk, n, right=of wp] (rules) {tikrinamos rūšies taisyklės};
  \node[man, n, right=of rules] (appr) {redaguota, patvirtinta};
  \draw[flow] (wp) -- (rules); \draw[flow] (rules) -- (appr);
  % --- joins
  \draw[opt] (reg.south) |- ($(gap.north)+(0,5mm)$) -- (gap.north);
  \draw[opt] (cit.south) -- ++(0,-4mm) -| ($(gap.north)+(-3mm,0)$);
  \draw[opt] (docs.west) -- ++(-3mm,0) |- (wp.west);
  % --- lanes
  \begin{scope}[on background layer]
    \node[lane, fit=(src)(reg)] (l1) {};
    \node[lane, fit=(docs)(cit)] (l2) {};
    \node[lane, fit=(app)(areg)] (l3) {};
    \node[lane, fit=(wp)(appr)] (l4) {};
  \end{scope}
  \node[lanelabel, anchor=south west] at (l1.north west) {Teisė --- kartą kiekvienam aktui};
  \node[lanelabel, anchor=south west] at (l2.north west) {Įmonė --- jos dokumentai};
  \node[lanelabel, anchor=south west] at (l3.north west) {Ataskaita};
  \node[lanelabel, anchor=south west] at (l4.north west) {Programos dokumentas};
  % --- the chat
  \node[ai, n, text width=30mm, right=6mm of appr, minimum height=11.5mm] (chat) {pokalbių asistentas savo įrankiais skaito visa tai ir pradeda ataskaitas bei programos dokumentus};
\end{tikzpicture}
\actorlegend
\caption{Sistema kaip keturios juostos. Geltoni langeliai --- modelio iškvietimai, mėlyni ---
kodo atliekamos patikros ar skaičiavimai, tamsūs --- žmogaus sprendimai. Taškinės linijos
rodo, ką ataskaita paima iš aukščiau esančių juostų: taikomų aktų pareigas ir tai, ką sako
įmonės dokumentai.}\label{fig:intro:map}
\end{figure}

Trys paveiksle matomos taisyklės galioja visoje sistemoje; jos yra
\cref{ch:control} tema:

\begin{enumerate}
  \item \textbf{Teisė skaitoma, o ne prisimenama.} Aktą skaitančiam modeliui pateikiamas
    nuostatos tekstas iš teisės aktų bazės ir užduodamas klausimas apie jį, o ataskaita
    įmonę vertina pagal taip perskaitytas pareigas. Kai struktūrizuotame atsakyme nurodoma
    nuostata, kodas ją suranda; kai jame cituojamas dokumentas, kodas suranda citatą
    dokumente. Išimtis --- paties pokalbio tekstas: jo niekas netikrina, o pokalbiui nurodyta
    kiekvieną skaičių, kurio jis neperskaitė teisės aktų bazėje, pažymėti kaip nepatikrintą
    (\cref{ch:chat}).
  \item \textbf{Modelis rengia projektą, kodas tikrina, sprendžia žmogus.} Balai, rizikos
    lygiai, plano tvarka ir prie išvados priskirta bauda apskaičiuojami, o ne paimami iš
    modelio. Sistema pati niekada nieko nepažymi patvirtintu: registro pareigą, peržiūrėtojo
    pakeistą išvados būseną, ataskaitą ir programos dokumentą kiekvieną
    patvirtina įvardytas žmogus, ir patvirtinimas užregistruojamas. Ataskaitos apimtis
    patvirtinama taip pat; vietoj to žmogus gali leisti ataskaitai vykti be sustojimo, ir
    tada ataskaitoje užfiksuojama, kad jos apimties niekas nepatvirtino. Pokalbių asistentas
    savo iniciatyva nieko netvirtina. Jis gali įrašyti tris sprendimus, visada naudotojo
    vardu: ataskaitos apimtį --- naudotojui pritarus ir už naudotojo priimtą kainą;
    peržiūrėtojo pakeistą išvados būseną --- naudotojo sprendimu; ir leidimo pažymėjimą
    patikrintu --- naudotojui pasakius, kad leidimas buvo patikrintas. Jis netvirtina nė vienos
    pareigos, programos dokumento ar ataskaitos
    (\cref{ch:chat}).
  \item \textbf{Kiekvienas iškvietimas užregistruojamas.} Kiekvienas atsakymą grąžinęs modelio
    iškvietimas registruojamas su jo žetonais, kaina ir, jei jį sukūrė instrukcija,
    instrukcijos versija (sąmoningai praleidžiamas tik paieškos užklausos vektorizavimas ---
    keli žetonai); kiekviena instrukcija saugoma versijomis, kurias administratorius gali
    skaityti, keisti ir grąžinti ankstesnę versiją; kiekvienoje ataskaitoje įrašyta, pagal
    kurias instrukcijų versijas ir kurias teisės aktų redakcijas ji parengta.
\end{enumerate}

\section{Kur dirba modeliai}\label{sec:intro:models}

Kalbos modelis iškviečiamas dvidešimt trijose vietose, kiekvienoje su savo saugoma
instrukcija, o kitos rūšies modelis paieškai paverčia tekstą vektoriais.
\Cref{tab:intro:paths} juos sugrupuoja; šio dokumento II dalyje jie aptariami po vieną, o
\cref{app:prompts} aprašo kiekvieną instrukciją.
Visi jie pasiekiami per vieną paslaugą ir vieną DI tinklų sąsają (\cref{ch:models}).

\begin{table}[H]
\caption{Visos vietos, kuriose iškviečiamas modelis, pagal paskirtį}\label{tab:intro:paths}
\begin{tabularx}{\textwidth}{@{}L{31mm}YL{19mm}@{}}
\toprule
\hd{Paskirtis} & \hd{Ko prašoma modelio} & \hd{Skyrius} \\
\midrule
Teisės skaitymas, kartą kiekvienam aktui &
  Užrašyti akto nuostatų pareigas (\prompt{corpus.extract_obligations}) ir standarto
  reikalavimus (\prompt{corpus.extract_requirements}); pasakyti, kam taikomos paties akto
  taikymo srities nuostatos (\prompt{corpus.scope}); suskirstyti jo sankcijų nuostatas į
  pakopas (\prompt{corpus.sanctions}); pažymėti antraštes standarto, kurio numeracija
  neįprasta (\prompt{corpus.standard_outline}); ir, kaip vertintojui, patikrinti kiekvieną
  išskirtą pareigą pagal nuostatą, kurią ji nurodo (\prompt{corpus.verify_obligations}). &
  \cref{ch:corpus,ch:obligations} \\
Įmonės dokumentų skaitymas &
  Sudaryti inventorių to, ką nustato dokumentas (\prompt{report.inventory}); ir, kaip
  vertintojui, pasakyti, ar straipsnis, kurį nurodo dokumentas, reglamentuoja tai, ką
  dokumentas teigia (\prompt{corpus.check_citations}). &
  \cref{ch:company} \\
Ataskaitos rengimas &
  Nuspręsti, ar aktas taikomas įmonei (\prompt{report.applicability}); įvertinti kiekvieną
  pareigą pagal dokumentus (\prompt{report.gap_analysis}); išskirti iš politikų įmonės
  vidaus reikalavimus ir juos įvertinti (\prompt{company.internal_requirements},
  \prompt{report.internal_gap}); suplanuoti veiksmus ir sujungti projektus
  (\prompt{report.action_plan}, \prompt{report.plan_merge}); parašyti santrauką valdybai
  (\prompt{report.summary}). &
  \cref{ch:report} \\
Programos dokumentų rengimas &
  Parengti atitikties politikos projektą ir išvardyti trūkstamus dokumentus
  (\prompt{programme.policy}, \prompt{programme.documents}); valdymo struktūrą ir atsakomybių
  matricą (\prompt{programme.structure}, \prompt{programme.raci}); rizikos kriterijus ir
  scenarijus (\prompt{programme.risk_criteria}, \prompt{programme.risk_scenarios}). &
  \cref{ch:programme} \\
Pokalbis &
  Palaikyti pokalbį ir pasirinkti, kuriuos įrankius iškviesti (\prompt{chat.assistant});
  suteikti pokalbiui pavadinimą (\prompt{chat.title}). &
  \cref{ch:chat} \\
Paieška &
  Vektorizavimo modelis paverčia vektoriais teisės aktų bazės nuostatas, pareigas, įmonės
  dokumentų ištraukas ir kiekvieną paieškos užklausą. &
  \cref{ch:models,ch:corpus} \\
\bottomrule
\end{tabularx}
\end{table}

Ne mažiau iškalbinga ir tai, kur modelis neiškviečiamas. Teisės atsisiuntimas ir jos
struktūros atpažinimas (išskyrus standarto antraštes, kurių neįstengia atpažinti taisyklės), citatos
suradimas, nuostatos paieška, išvados susiejimas su jos sankcija, rizikos balo
apskaičiavimas, plano išrikiavimas, veiksmų registro keitimas ir kiekvienas patvirtinimas yra
kodo arba žmonių darbas. Ypač veiksmų registre --- toje sistemos dalyje, kuri pasako, ar
spraga iš tikrųjų panaikinta --- nėra nė vieno modelio iškvietimo (\cref{ch:register}).

\section{Sistema skaičiais}\label{sec:intro:numbers}

Toliau pateikti skaičiai buvo nuskaityti įdiegtoje sistemoje 2026~m. lapkričio 13~d.; nuo tada
pasikeitę du skaičiai --- instrukcijų versijų ir sankcijų pakopų, dar kartą perskaitytų teisės
aktų bazės kopijoje nauja jų instrukcijos versija, --- pateikiami tokie, kokie yra nuo
2026~m. lapkričio 19~d. Jie pateikiami tam, kad
skaitytojas susidarytų mastelio įspūdį; skyriuose paaiškinta, kaip kiekvienas gautas.

\begin{table}[H]
\caption{Sistema 2026 m. lapkričio 13 d.}\label{tab:intro:numbers}
\begin{tabularx}{\textwidth}{@{}L{44mm}Y@{}}
\toprule
\hd{Kas} & \hd{Kiek} \\
\midrule
Teisės aktų bazė & 78 aktai kataloge, iš jų 68 įjungti; 26\,983 nuostatos, saugomos kaip
  vienetai galiojančiomis redakcijomis; 4\,059 ištraukos, vektorizuotos paieškai
  (\cref{app:acts}) \\
Pareigų registras & 5\,960 pareigų, išskirtų iš 66 aktų iš 68 (du aktai pareigų nenustato): 5\,938
  patvirtintos, 10 atmesta, 12 laukia žmogaus; 302 sankcijų pakopos, perskaitytos iš 27 aktų \\
Instrukcijos & 23, nuo 2026~m. lapkričio 20~d. iš viso 91 versija (iki lapkričio 16~d. --- 85); pokalbio
  instrukcija yra septynioliktos versijos,
  ataskaitos santraukos --- dvyliktos \\
Viena visa ataskaita & Parengta 2026~m. lapkričio 13~d. be apimties peržiūros kurso pavyzdinei įmonei,
  kriptoturto paslaugų teikėjui su 26 politikomis: įvertinti 26 aktai (vienas jų --- kaip
  pasirengimas dar būsimai datai), 2\,151 išvada pagal teisę ir 475 pagal pačios įmonės
  politikas, 51 veiksmo veiksmų planas, apimantis visas 1\,633 neišspręstas išvadas;
  272 modelio iškvietimai, 1 val. 38 min., 13,20~\$ (\cref{ch:report}) \\
Kodas & Apie 33\,600 eilučių API dalyje, 10\,700 asistento paslaugoje ir 46\,800 konsolėje, šalia
  jų --- apie 94\,000 testų eilučių \\
Testai & 405 API dalyje, 509 asistento paslaugoje, 5\,180 konsolėje (\cref{ch:testing}); šeši
  pažymėtų atvejų rinkiniai, kuriais šešios instrukcijos matuojamos su tikruoju modeliu;
  visas žymas 2026 m. lapkričio 12 d. sukūrė sistemai kurti naudotas DI programavimo
  asistentas, ir nė viena dar nepatvirtinta žmogaus (\cref{ch:control}) \\
\bottomrule
\end{tabularx}
\end{table}

\section{Kursas, kuriam sukurta sistema}\label{sec:intro:course}

\product{} yra Vilniaus universiteto Kauno fakulteto FinTech magistrantūros programos kurso
\enquote{Finansų teisė ir finansų įstaigų veiklos atitikties valdymas} projektas; kursas dėstomas
2026 m. rudenį. Jis moko metodo: klasifikuoti įstaigą, sukurti jos atitikties politiką ir
valdymą, įvertinti jos teisinę riziką, stebėti ir teikti ataskaitas. Jame pateikiamos trys
kurso užduotys apie pavyzdinę organizaciją \demo{} --- išgalvotą kriptoturto paslaugų teikėją
su 26 politikomis, --- ir aprašomas technologo, kuriančio įrankius atitikties darbui,
vaidmuo: su schemomis susietos instrukcijos, rastais šaltiniais pagrįsti atsakymai, saugomi
pažymėtų atvejų rinkiniai, regresijos matavimas, veiksmų žurnalai. Sistema yra bandymas šį
metodą įgyvendinti programine įranga, o pavyzdinė organizacija yra ta įmonė, kuriai buvo
parengtos šio dokumento ataskaitos ir atlikti kiti su įmone susiję matavimai.
\Cref{ch:course} dalis po dalies susieja kurso užduotis ir sąvokas su sistema.

\section{Ko sistema neteigia}\label{sec:intro:limits}

\begin{itemize}
  \item \textbf{Tai nėra teisinė konsultacija.} Ataskaita yra struktūrizuotas pirminis
    perskaitymas, kurį turi peržiūrėti atitikties pareigūnas: jos santraukoje yra skirsnis
    apie jos pačios ribotumus, ir joje užfiksuota, kiek jos išvadų remiasi patvirtintomis
    registro pareigomis.
  \item \textbf{Ji vertina tai, kas užrašyta.} Išvada pasako, ar įmonės dokumentai nustato
    tai, ko reikalauja pareiga, --- t. y. įmonės nustatytos tvarkos sandarą. Ar įmonė praktiškai daro tai,
    ką sako jos politika, dokumentai parodyti negali; išvadoje, kurios būsena
    \status{met}, yra pastaba tai patikrinti praktikoje, kai tai svarbu.
  \item \textbf{Ji skaito tekstą.} Nuskenuotame dokumente be teksto sluoksnio jai nėra ko
    skaityti: optinio simbolių atpažinimo nėra.
  \item \textbf{Jos teisė yra jos teisės aktų bazės teisė.} Ji apima kataloge išvardytą
    Europos Sąjungos ir Lietuvos finansų reguliavimą (\cref{app:acts}) ir bet kurį aktą,
    kurį prideda administratorius. Teisės aktų bazei nepriklausantis aktas jai neegzistuoja,
    ir ji tai pasako, kai dokumentas tokį aktą nurodo.
  \item \textbf{Jos matavimai yra stebėjimai.} Šiame dokumente pateikti tikslumo rodikliai,
    kainos ir trukmės buvo išmatuoti vienos pavyzdinės įmonės duomenimis arba bendroje teisės aktų bazėje
    nurodytomis dienomis. Jie rodo, kaip sistema elgėsi, o ne kaip ji turi elgtis kitur.
\end{itemize}

Kiekvienos dalies ribotumai nurodyti jos skyriuje ir surinkti \cref{ch:control}.
